% Generado por sistema/backend ReporteTex.java -- no editar a mano.
% Datos: dataset de prueba (respuestas sinteticas)
\textit{Paso 1. Contar con la función de probabilidad} $p(x) = f_i/n$, a partir de las frecuencias observadas en la encuesta. \textit{Paso 2. Definir la acumulada} $F(x_i) = F(x_{i-1}) + p(x_i)$. \textit{Pasos 3 y 4. Igualar $F(x) = R$ e invertir:} a cada categoría le corresponde el rango $[F(x_{i-1}),\,F(x_i))$, y $x = F^{-1}(R)$ es la categoría cuyo rango contiene a $R$.

\begin{mitabla}{Transformada Inversa de la situación laboral, escenario base ($n=100$) \datosprueba}
\begin{tabular}{lcccc}
\toprule
$x$ & $f_i$ & $p(x)$ & $F(x)$ & Rango de $R$ \\
\midrule
Trabaja en una organización & 68 & $68/100 = 0.680$ & 0.680 & $[0.000,\ 0.680)$ \\
Emprendimiento propio & 19 & $19/100 = 0.190$ & 0.870 & $[0.680,\ 0.870)$ \\
No trabaja & 13 & $13/100 = 0.130$ & 1.000 & $[0.870,\ 1.000)$ \\
\bottomrule
\end{tabular}
\end{mitabla}

\begin{mitabla}{Transformada Inversa del tiempo al primer empleo, escenario base ($n=68$) \datosprueba}
\begin{tabular}{lccccc}
\toprule
$x$ & $f_i$ & Meses & $p(x)$ & $F(x)$ & Rango de $R$ \\
\midrule
Ya trabajaba & 7 & 0 & $7/68 = 0.103$ & 0.103 & $[0.000,\ 0.103)$ \\
$<$ 1 mes & 12 & 0.5 & $12/68 = 0.176$ & 0.279 & $[0.103,\ 0.279)$ \\
1–4 meses & 15 & 2.5 & $15/68 = 0.221$ & 0.500 & $[0.279,\ 0.500)$ \\
4–8 meses & 11 & 6 & $11/68 = 0.162$ & 0.662 & $[0.500,\ 0.662)$ \\
8–12 meses & 12 & 10 & $12/68 = 0.176$ & 0.838 & $[0.662,\ 0.838)$ \\
$>$ 12 meses & 11 & 12 & $11/68 = 0.162$ & 1.000 & $[0.838,\ 1.000)$ \\
\bottomrule
\end{tabular}
\end{mitabla}

\textit{Escenario modificado.} Se traslada el 50\% de la probabilidad de ``No trabaja'' a ``Trabaja en una organización'':
\[
p^{\prime}(\text{No trabaja}) = 0.130 \times (1 - 0.50) = 0.0650, \qquad p^{\prime}(\text{Trabaja}) = 0.680 + 0.130 \times 0.50 = 0.7450
\]
y el 30\% de la probabilidad de los tiempos mayores a 4 meses a ``1--4 meses'':
\[
p^{\prime}(\text{1--4 meses}) = 0.221 + (0.162 + 0.176 + 0.162) \times 0.30 = 0.3706
\]
y cada uno de esos tres rangos queda multiplicado por $1 - 0.30$.

\begin{mitabla}{Transformada Inversa de la situación laboral, escenario modificado \datosprueba}
\begin{tabular}{lccc}
\toprule
$x$ & $p(x)$ & $F(x)$ & Rango de $R$ \\
\midrule
Trabaja en una organización & 0.745 & 0.745 & $[0.000,\ 0.745)$ \\
Emprendimiento propio & 0.190 & 0.935 & $[0.745,\ 0.935)$ \\
No trabaja & 0.065 & 1.000 & $[0.935,\ 1.000)$ \\
\bottomrule
\end{tabular}
\end{mitabla}

\begin{mitabla}{Transformada Inversa del tiempo al primer empleo, escenario modificado \datosprueba}
\begin{tabular}{lcccc}
\toprule
$x$ & Meses & $p(x)$ & $F(x)$ & Rango de $R$ \\
\midrule
Ya trabajaba & 0 & 0.103 & 0.103 & $[0.000,\ 0.103)$ \\
$<$ 1 mes & 0.5 & 0.176 & 0.279 & $[0.103,\ 0.279)$ \\
1–4 meses & 2.5 & 0.371 & 0.650 & $[0.279,\ 0.650)$ \\
4–8 meses & 6 & 0.113 & 0.763 & $[0.650,\ 0.763)$ \\
8–12 meses & 10 & 0.124 & 0.887 & $[0.763,\ 0.887)$ \\
$>$ 12 meses & 12 & 0.113 & 1.000 & $[0.887,\ 1.000)$ \\
\bottomrule
\end{tabular}
\end{mitabla}

